\section{The pure conductor-four four-residue window}
\label{sec:pure-conductor4-quarter-actual-window}

\paragraph{Status and exact dependency.}
This conditional-interface proof uses the literal original-coefficient
strict theorem AC-G162.  It retains every omission in AC-G159-window
and enlarges only the weight progression for the already present
pure quadratic character $\delta_{-4}$, with trivial cubic part.
Its separately proved class and unit inputs are not substituted for
the full strict theorem.  One local comparison cost covers the entire union.
The remaining all-prime estimate and mixed realization stay open.

Put $p=269$, $K=\mathbb Q(\omega)$, and $t=\sqrt[3]p>0$;
let $\chi:G_K\to\mu_3$ describe $K(t)/K$.
Let $\phi$ be the primitive odd quartic character modulo sixteen
with $\phi(13)=i$ and kernel $\{1,7\}$, and put
$\delta_d=\chi_d|_{G_K}$ for $d=-4,8,-8$.
For the retained datum, $0<\varepsilon<1/256$, allowed page and
prescribed pair $q,d$, keep the actual objects
\[
 \begin{aligned}
 F_u&=\mathbb Q(\zeta_{2^{h(u)+2}},\zeta_7,
                           \sqrt{-q},\sqrt{-d},\omega),\\
 H_u&=F_u(t),\quad G_u=\operatorname{Gal}(H_u/K),\quad D_u=|G_u|,\\
 u^{2+8\varepsilon}&\leq2^{h(u)}\leq2u^{2+8\varepsilon},
                         \qquad T=\lfloor u^{399}\rfloor .
 \end{aligned}
 \tag{P4QW.1}
\]
The prescribed pair's letter $d$ is independent of the subscript
of $\delta_d$.
All original coefficient rings $O_\ell$, fields $E_\ell$,
normalizing ranks $b_\ell$, full local conditions and derived
indices retain (FC.1)--(FC.2), (EU.2) and (TC.1).
In particular the coefficient-prime range is unrestricted.

For $u\geq2$ the new omitted union is
\[
 \begin{aligned}
 Z_{16}^+&=\{(p,\chi^{-j}\phi^{-\sigma},e):j=1,2,\ \sigma=\pm1,\\
 &\hspace{2em}1\leq e<T,\ e\equiv-1\pmod{67}\},\\
 Z_4^+&=\{(p,\chi^{-j}\delta_{-4},e):j=1,2,\ 1\leq e<T,
                                      \ e\equiv-1\pmod{67}\},\\
 Z_8^+&=\{(p,\chi^{-j}\delta_d,e):j=1,2,\ d=8,-8,\\
 &\hspace{2em}1\leq e<T,\ e\equiv-1\pmod{67}\},\\
 Z_{\rm cub}^{(4)}&=\{(p,\chi^j,e):j=0,1,2,\ 1\leq e<T,\\
 &\hspace{2em}e\equiv0\pmod{268}\ \text{or}\ e\equiv-1\pmod{67}\},\\
 Z_3&=\{(3,1,e):1\leq e<T\},\\
 Z_{4}^{0,+}&=\{(p,\delta_{-4},e):1\leq e<T,
                                      \ e\equiv-1\pmod{67}\},\\
 Z_{8}^{0}&=\{(p,\delta_d,e):d=8,-8,\ 1\leq e<T,
                                      \ e\equiv-1\pmod{134}\},\\
 Z^{0,4q}&=Z_{16}^+\sqcup Z_4^+\sqcup Z_8^+
         \sqcup Z_{\rm cub}^{(4)}\sqcup Z_3\sqcup Z_4^{0,+}\sqcup Z_8^0.
 \end{aligned}
 \tag{P4QW.2}
\]
Only the pure conductor-four progression enlarges its AC-G159
predecessor; the two pure quadratic-eight progressions remain unchanged.
The added subset is
\[
 Z_{4,\mathrm{new}}^0=
 \{(p,\delta_{-4},e):1\leq e<T,
                                  \ e\equiv66,200\pmod{268}\}.
 \tag{P4QW.3}
\]
Let $I_{399}^{\rm estpure4quarter}$ be (TC.1) with precisely
$Z^{0,4q}$ omitted; the previous remaining sum is
$I_{399}^{\rm est8quarter}$.

\begin{theorem}[AC-G163-window: the pure conductor-four four-residue reduction]
\label{AC-G163-window}
Every label in (P4QW.2) is an actual character of the prescribed
host, and its strict group on the original coefficient module
vanishes.  For positive integer $T$ the exact disjoint counts are
\[
 \begin{aligned}
 |Z_{16}^+|&=4\lfloor T/67\rfloor,& |Z_4^+|&=2\lfloor T/67\rfloor,\\
 |Z_8^+|&=4\lfloor T/67\rfloor,& |Z_3|&=T-1,\\
 |Z_{\rm cub}^{(4)}|&=3\bigl(\lfloor(T-1)/268\rfloor
                                      +\lfloor T/67\rfloor\bigr),\\
 |Z_4^{0,+}|&=\lfloor T/67\rfloor,\qquad
 |Z_8^0|=2\lfloor T/134\rfloor,\\
 |Z_{4,\mathrm{new}}^0|&=\lfloor(T+67)/134\rfloor.
 \end{aligned}
 \tag{P4QW.4}
\]
Let $C_u$ be the actual normalized positive-index cost over
this entire union.  Then
\[
 \begin{aligned}
 I_{399}&=I_{399}^{\rm estpure4quarter}+C_u,\qquad
 I_{399}^{\rm estpure4quarter}\leq I_{399}^{\rm est8quarter}\leq I_{399},\\
 0\leq C_u&\leq8D_u^2T^2\log(8TD_u)
                   +\tfrac12D_u\mathcal R_uT^2+2D_ur_uc_0T\\
 &=O_\varepsilon\bigl(u^{802+16\varepsilon}(\log u)^2\bigr).
 \end{aligned}
 \tag{P4QW.5}
\]
The first line is an identity and monotonicity statement for
nonnegative extended sums.
One full (FL.4) error covers all old and new omitted labels;
no individual positive index is asserted zero.

The original common-window premise (TC.3), or its AC-G159 version,
may use $I_{399}^{\rm estpure4quarter}$ with every simultaneous field,
coefficient, finite-character and weight quantifier unchanged.
Either older premise implies this one with unchanged parameters.
Conversely the new positive-relative-slack premise implies the
full original premise on the same $W,J$, datum, page and pair,
after only increasing its initial threshold and halving slack.
The actual $R.t$, every other term and all clauses of the retained
mixed input $\mathcal R_{\rm cub}$ remain required and unproved.
\end{theorem}

\begin{proof}
\proofstep{1. Actual labels and the exact original modules.}
The host contains $\mathbb Q(\zeta_8)$, so $\delta_{-4}$ is
an actual character already present in AC-G159, with trivial cubic part.
The old distinctness proof is unchanged: $F_u\cap K(t)=K$
separates the cubic factor, and the distinct two-power characters
remain distinct on $K$.  A rational ratio ramified only at two
cannot equal $\eta_K$, ramified at three; order-four and
order-two restrictions also remain distinct.
The pure $\delta_8,\delta_{-8}$ labels retain the smaller
progression $e\equiv-1\pmod{134}$.
No new finite character at $269$ is introduced.

For each added term the original coefficient is literally
\[
 \begin{aligned}
 A_e&=(O_p\otimes_{\mathbb Z_p}\mathbb Q_p/\mathbb Z_p)
                                                   (\delta_{-4})(-e)\\
    &=(E_p/O_p)(\delta_{-4}c_p^{-e}).
 \end{aligned}
 \tag{P4QW.6}
\]
The original $O_p$ contains the values $\{1,-1\}$.
Apply (P4Q.2) of AC-G162 with $m=0$ and $\kappa=1$.
It proves exactly the strict zero of (P4QW.6) at
$e\equiv66,200\pmod{268}$, with arbitrary absolute
ramification of this same $O_p$.
The quartic character of conductor $269$ occurs only in the
residual Kummer cover used to prove that theorem.
This is the literal cubic-index-zero theorem; the nontrivial-cubic
AC-G152 theorem is not used to infer it.
Every old strict zero, including the other two pure quadratic rows
and the coefficient-three theorem without averaging by three,
is retained from AC-G159.

\proofstep{2. Full local conditions and exact union counts.}
The source theorem imposes zero in the full
$H^1(G_{K_v},A_e)$ at each $v\mid p$ and zero in
$H^1(I_v,A_e)$ at every other finite place.
These are the original (FC.2) conditions.
The coefficient is unramified outside $2,3,p$.
At an added away-$p$ prime of $\Sigma_p$, its inertia action is
trivial, so a coboundary cocycle restriction is zero.
The cocycle identity gives zero on every conjugate inertia
subgroup and their generated closed normal subgroup.
The coefficient action and cocycle descend to the smaller
ramification set; conversely inflation is injective and preserves
all remaining local conditions.  In particular full local zero
at $p$ is retained.
Thus the theorem applies to the actual original strict modules
in (FC.2), including the host's enlarged ramification set.

The new weights are $e=66+134n$, $n\geq0$.
For $1\leq e<T$ their number is $\lfloor(T+67)/134\rfloor$.
The old pure conductor-four weights have the form $133+134n$.
The two disjoint progressions combine to $e=67n-1$, $n\geq1$,
so this one finite label now has $\lfloor T/67\rfloor$ weights.
The two other pure quadratic labels still contribute
$2\lfloor T/134\rfloor$ weights in total.
These statements prove all changed counts in (P4QW.4);
every other count is retained literally from AC-G159.
The new strict terms disappear from (FC.15), but their derived
positive-index costs are still included in $C_u$.

\proofstep{3. One total-union comparison cost.}
Use the original dual complex
$C_{\ell,\eta,e}=R\Gamma(G_{K,\Sigma_\ell},
O_\ell(\eta^{-1})(e+1))$, with the codifferent normalization
of (FL.2).  For every removed label, (EU.2) and strict zero give
\[
 \begin{aligned}
 \frac{\log\#k_{O_\ell}}{b_\ell}\max(0,i_{\ell,\eta,e})
 &\leq\frac{\log\#H^2(C_{\ell,\eta,e})}{b_\ell}\\
 &=\frac{|\log\#H^2(C_{\ell,\eta,e})
                  -\log\#\mathcal S_{\ell,\eta}(e)|}{b_\ell}.
 \end{aligned}
 \tag{P4QW.7}
\]
The positive index and the entire original normalizing rank
$b_\ell$ are retained.  At each weight the whole union is one
subsum of the nonnegative all-prime absolute error (FL.4),
so its cost is at most
\[
                  16eD_u^2\log(8eD_u)+D_ue\mathcal R_u+2D_ur_uc_0.
 \tag{P4QW.8}
\]
Summing once over $1\leq e<T$, using $\sum e\leq T^2/2$
and $\log(8eD_u)\leq\log(8TD_u)$, proves (P4QW.5).
No extra full error is charged for the new family.
Nonnegative additivity proves its extended-sum identity and
monotonicity without subtracting infinite quantities.

The unchanged host estimates are $D_u\leq288u^{2+8\varepsilon}$,
$r_u\leq12$ and
$\mathcal R_u=O_\varepsilon(u^{\varepsilon/5}(\log u)^2)$.
Together with $T\leq u^{399}$ they give largest exponents
$802+16\varepsilon$ and $800+41\varepsilon/5$.
Hence $C_u=o(\exp(u^\varepsilon/2))$, uniformly over the
prescribed page and pair, for each fixed positive $\varepsilon$.

\proofstep{4. The original common window and slack.}
Fix the retained datum, page and $0<\varepsilon<1/256$.
Suppose the new premise supplies $J\geq1$, $X_0$ and
$s\in(0,1)$ independent of $X,q,d,u,\ell,\eta,e$, such that
for every $X\geq X_0$ and prescribed
$\operatorname{WinPair}(\varepsilon,X^{1/J},q,d)$ there is
one $W$ with $X^{1/J}\leq W$, $W^2\leq X$ and
\[
 I_{399}^{\rm estpure4quarter}(H_u,u)\leq(1-s)\exp(u^\varepsilon/2)
                                      \quad(W\leq u\leq W^2).
 \tag{P4QW.9}
\]
Choose $U_0\geq2$ making (P4QW.5) at most
$(s/2)\exp(u^\varepsilon/2)$ for every $u\geq U_0$.
Only the retained datum, page, $\varepsilon$ and slack enter
this uniform choice.  Replace $X_0$ by $\max(X_0,U_0^J)$.
The very same $W$ then gives
\[
 I_{399}(H_u,u)\leq(1-s/2)\exp(u^\varepsilon/2)
                                      \quad(W\leq u\leq W^2).
 \tag{P4QW.10}
\]
For a fixed finite simultaneous family use the maximum of its
uniform error constants while keeping this one $W$.
There is no intersection of separately chosen windows.
Conversely either old premise implies (P4QW.9) by nonnegativity,
with unchanged parameters.  Uniform little-$o$ versions agree.
Every host, original coefficient, actual source $R.t$, finite
comparison cost and mixed realization clause remains as before.
\end{proof}

\paragraph{Exact remaining scope.}
For $T\geq67$ the added set (P4QW.3) is nonempty and disjoint
from the old omitted union, so the sufficient premise covers a
proper smaller set of actual terms.
This does not assert a strict numerical decrease of the sum.
Every unlisted coefficient prime, finite character and weight,
the all-prime numerical bound, mixed theta/Selmer identification
and analytic realization, and H04.4 remain open.
Only the previously audited SGA~1, Expos\'e~V, \S\S4 and 7
finite-cover and original rank-one Artin--Tate dictionary is used.
No new motive, comparison theorem, maximal vanishing classification
or Lean execution is asserted.
